% Transcription — fonds Alexandre Grothendieck, shelfmark 90, batch 1 (pages 1-20).
% Established from the facsimile archives/batches/90/batch-01.pdf (a copy of
% 90.pdf fetched through the project's relay,
% https://grothendieck-relay.onrender.com/source/90.pdf; PDF page k =
% archivists' page k, no cover sheet), per the skill transcribe-grothendieck.
% The folder is twenty-four pages; this is the first of two batches (pages
% 21-24 are batch 2, transcribed separately; its notation was read before this
% file was finished and is followed here: \underline{\Gamma}^{\cdot}, A^{\cdot},
% \widetilde{M}, O_X-G-Module).
% Pass: Opus 5.5 (claude-opus-5-5), 2026-09-26 — first pass, unchecked against the pages by a human.
%
% Pages skipped: none.
%
% Pages summarised: every even page (2, 4, 6, 8, 10, 12, 14, 16, 18, 20).
% The notes are written on the rectos of leaves whose other side is typescript
% of EGA IV, §§ 18.6-18.9, scanned upside down (read turned). Three runs:
% pp. 4 and 2, « EGA IV 18.9. (Rabiots) », Th. 18.9.7 and Cors. 18.9.8-9 with
% the end of a proof; pp. 6 and 14, a retyped Cor. 18.9.5 paginated « -3- »,
% « -4- »; pp. 8, 10, 12, 16, 18, 20, leaves paginated by hand IV-965, 964,
% 963(?), 962, 953, 952(?) (§18.6 henselisation, §18.8 strict henselisation,
% §18.9 formal fibres of henselian rings; batch 2's pp. 22 and 24 are IV-951
% and IV-950 of the same run). The typed text is published and not re-set:
% one French \note per page, with the ink interventions listed. As in batch 2,
% the hand of the ink corrections is not established (they may be the
% redactor's); the notes say so. Page 14 carries only a completed letter, and
% is kept with a one-line note for the sake of the -3-/-4- pair.
%
% Pages transcribed from his hand: all odd pages, 1-19, blue ink (pp. 1, 3,
% 17, 19 darker; 5-15 paler, thinner pen). Fast drafting: formulas carry,
% much connecting prose does not. Reading order: 1, 3 (Lemmas 1-2 on a
% Dedekind base), 5 (two corollaries), 7 -> 9 -> 11 (a Proposition on
% A-G-modules through the dual algebras U_i, and its sequel), 13 -> 15 (a
% cancelled Proposition, then a general statement 1)-4) over S), 17 -> 19 ->
% batch 2 p. 21 -> p. 23 (the clean restatement for a flat affine G over S
% acting on X with an ample L, and its proof through graded modules).
% His pagination: none.
%
% What the batch is about. Not K-groups: nothing on these pages names K(X),
% K_G or Riemann-Roch, despite the folder title « KG (X) ». What they hold is
% the groundwork an equivariant K-theory needs — that equivariant coherent
% sheaves are quotients of equivariant locally free ones. P. 1: Lemma 1, for
% A Dedekind and M a flat A-module, conditions equivalent to M being the
% filtered union of saturated finite-type submodules (N~ = inverse image of
% the torsion of M/N; M = lim M_i with M_i saturated; sufficient: linear forms
% separate M, M -> A^I; necessary too for countably generated M). P. 3: Lemma
% 2, for G affine over A with algebra B satisfying Lemma 1, every finite-type
% M in B lies in a saturated N with Delta N in N (x) N (G x G acting by
% (g,h)f(s) = f(g^-1 s h)). P. 5: corollaries — G of finite type over a
% Dedekind A has a faithful linear representation; over a field every affine
% group is a strict projective limit of affine groups of finite type. Pp.
% 7-11: if B = lim B_i with B/B_i flat, Delta B_i in B_i (x) B_i, B_i
% projective of finite type, then U_i = dual of B_i is an associative unital
% algebra, the category of A-G-modules of finite type is the inductive limit
% of the categories of U_i-modules of finite type, and every G-module of
% finite type is a quotient of a projective one; a Proposition for G flat
% affine over a Dedekind ring and a scheme S' over S on which quasi-coherent
% sheaves of finite type are quotients of locally free ones. Pp. 13-15, 17-19:
% the general form over S — for X over S with G acting, E a quasi-coherent
% O_X-G-Module: 1) the smallest G-stable quasi-coherent F' containing a given
% F exists, commutes with flat base change, and is of finite type if F is;
% 2) E is the inductive limit of its G-stable quasi-coherent submodules of
% finite type; 3) a surjection F (x) B_i^v -> E from a G-module locally free
% if F is; 4) under an ample L (or X noetherian regular) every O_X-G-Module
% of finite type is a quotient of a locally free one of the form f*(H)(-n).
% The proof (p. 19, continued on batch 2's pp. 21, 23) goes through
% Gamma_*(E) = coprod_n f_*(E(n)), a graded module over A^. = Gamma^.(O_X) on
% which G acts. Batch 2's reading of its own pages fits this: p. 19 ends « On
% a un foncteur (exact) » and p. 21 opens « M -> M~ en sens inverse ».
%
% Notation fixed for the batch: \widetilde{N} for his tilde (saturation) on
% pp. 1-5; \Delta for the comultiplication; \check{B}_i for the dual of B_i;
% \mathcal{B}, \mathcal{B}_i on pp. 13, 17 where the B is drawn as a script
% sheaf letter; \underline{O}_X; \underline{\Gamma}^{\cdot} as in batch 2.
% « CorMain », twice (pp. 5, 9), is kept as \uncertain{Cor. Main}; it may be
% « Corollaire ».
%
% Readings to check first: the whole of conditions a)-c) of Lemma 1 (p. 1,
% heavily reworked); « D.— » on p. 1; the middle of the proof of Lemma 2
% (p. 3); the sentence « Alors ... de G » on p. 5; the end of the cancelled
% Proposition on p. 13 (« ... fidèle et plat »?); the numbers of the typescript
% leaves on pp. 12 and 20; the margin notes on pp. 8 and 20.
%
% Apparatus count: 183 \ill{}, 37 \uncertain{}.

\documentclass[11pt,a4paper]{article}
% Le préambule commun à toutes les transcriptions du fonds Grothendieck.
%
% Il tient en une idée : **l'incertitude se compose, elle ne se résout pas.**
% Une transcription qui lisse un mot douteux a détruit la seule chose pour
% laquelle on la faisait — un lecteur doit pouvoir savoir, à chaque endroit,
% ce qui était sur le feuillet et ce qui a été ajouté. Les six macros qui
% suivent sont donc l'appareil critique, et non de la décoration.
%
% Les mêmes commandes sont reconnues par `scripts/render.mjs`, qui produit la
% vue de lecture du site. Une macro ajoutée ici doit l'être aussi là-bas,
% faute de quoi le rendu échouera bruyamment — ce qui est voulu : un rendu
% silencieusement faux est pire qu'un rendu absent.

\ProvidesPackage{grothendieck}[2026/08/08 Transcriptions du fonds Grothendieck]

\usepackage{amsmath,amssymb,amsthm}
\usepackage{mathrsfs}
\usepackage{stmaryrd}
% Les diagrammes commutatifs sont partout dans ce fonds : c'est la langue
% même de la géométrie algébrique de Grothendieck, et une transcription qui
% les rendrait en prose ne transcrirait rien.
\usepackage{tikz-cd}
% Français par défaut — c'est la langue des feuillets — mais l'anglais est
% chargé aussi : la traduction et le résumé sont des documents anglais, et la
% typographie française (espaces avant les deux-points) y serait une faute.
% Ces documents basculent par \selectlanguage{english} en tête de corps.
\usepackage[english,french]{babel}
\usepackage{fontspec}
\usepackage{geometry}
\geometry{a4paper,margin=25mm,marginparwidth=28mm}
\usepackage{marginnote}
\usepackage{xcolor}
% ulem, not soul. `soul`'s \st and \ul are built on a character-by-character
% scan that XeLaTeX's font machinery breaks: the document fails with an
% undefined control sequence rather than degrading. ulem does the same two jobs
% and is Unicode-engine safe. `normalem` stops it hijacking \emph, which the
% transcriptions use for Grothendieck's own emphasis.
\usepackage[normalem]{ulem}
\usepackage{hyperref}
\hypersetup{colorlinks=true,linkcolor=black,urlcolor=black,pdfborder={0 0 0}}

% --- Métadonnées du lot -----------------------------------------------------
% Reprises telles quelles par le script de rendu, qui les affiche en tête de
% la vue de lecture. Elles fixent ce que le fichier prétend couvrir : sans
% elles, une transcription flotte sans point d'attache dans un fonds de seize
% mille feuillets.

\newcommand{\folder}[1]{\gdef\grothFolder{#1}}
\newcommand{\batch}[1]{\gdef\grothBatch{#1}}
\newcommand{\leaves}[2]{\gdef\grothFirst{#1}\gdef\grothLast{#2}}
\newcommand{\foldertitle}[1]{\gdef\grothFoldertitle{#1}}
\gdef\grothFolder{?}\gdef\grothBatch{?}\gdef\grothFirst{?}\gdef\grothLast{?}\gdef\grothFoldertitle{}

% --- Le résumé --------------------------------------------------------------
%
% Il ouvre la lecture modernisée, sous le titre « Résumé », et il est composé
% en petit. Ce n'est pas un ornement : c'est le seul passage du document qui
% ne soit pas des mathématiques, et un lecteur doit pouvoir le reconnaître
% avant de l'avoir lu — puis le sauter s'il connaît déjà le sujet, ou s'y
% arrêter s'il ne le connaît pas. Le filet qui le suit marque la reprise du
% corps.

\newenvironment{resume}
  {\small\setlength{\parskip}{0.45em}}
  {\par\medskip\hrule\medskip}

% --- Mots-clés ---------------------------------------------------------------
%
% En anglais, en clôture du résumé de la lecture modernisée : le vocabulaire
% moderne sous lequel chercher ce que les feuillets construisent. Le manifeste
% du site (`npm run manifest`) les extrait pour étiqueter la cote entière —
% cette ligne est la source unique des tags, il n'y a pas de fichier à côté.
\newcommand{\keywords}[1]{\par\smallskip\noindent
  {\footnotesize\textsc{Keywords} --- \textit{#1}}\par}

% --- L'appareil critique ----------------------------------------------------

% Le numéro de feuillet, en marge. C'est l'ancre qui permet de régler un
% désaccord sur une formule en regardant *un* feuillet plutôt qu'une chemise
% de six cents. Le site s'en sert aussi pour tourner les pages du fac-similé
% quand on fait défiler la transcription.
\newcommand{\leaf}[1]{%
  \marginnote{\footnotesize\color{gray!70}\textsf{#1}}%
  \label{leaf:#1}\ignorespaces}

% Illisible : on ne devine pas. Un mot inventé qui se lit comme les autres est
% le pire résultat possible.
\newcommand{\ill}{\textcolor{red!60!black}{[\,\ldots\,]}}

% Lecture douteuse : le mot est donné, et signalé comme incertain.
\newcommand{\uncertain}[1]{\uline{#1}}

% Ajout de l'éditeur — une lettre manquante, un mot sous-entendu.
\newcommand{\add}[1]{\textcolor{blue!50!black}{[#1]}}

% Biffé par Grothendieck lui-même. Ce qu'il a rayé fait partie du document :
% une rature dit souvent où la pensée a changé de direction.
\newcommand{\struck}[1]{\sout{#1}}

% Note du transcripteur, sur l'état matériel du feuillet.
\newcommand{\note}[1]{\par\smallskip{\small\color{gray!60!black}\textit{[#1]}}\par\smallskip}

% Note marginale de Grothendieck, distincte du corps du texte.
\newcommand{\marginal}[1]{\par\smallskip\noindent\hspace{1em}%
  {\small\color{gray!60!black}#1}\par\smallskip}

% --- Titre ------------------------------------------------------------------

\AtBeginDocument{%
  \begin{center}
    {\large\bfseries Fonds Alexandre Grothendieck --- cote n\textsuperscript{o}~\grothFolder}\\[2pt]
    {\itshape\grothFoldertitle}\\[4pt]
    {\small feuillets \grothFirst--\grothLast\ (lot \grothBatch)}\\[2pt]
    {\footnotesize\color{gray!60!black}Transcription établie d'après le fac-similé numérisé
     par l'Université de Montpellier.\\
     Les lectures incertaines sont soulignées, les illisibles notées [\,\ldots\,],
     les ajouts entre crochets.}
  \end{center}
  \vspace{1em}}

\endinput


\folder{90}
\batch{1}
\pages{1}{20}
\dating{[à partir de 1964]}
\watermark{Édition de démonstration}
\foldertitle{KG (X) : notes manuscrites (s.d.)}

\begin{document}

\page{1}

\struck{$f(gh) = \sum f_i(g)\,\varphi_i(h)$, $M$ \ill{}}
\note{les deux premières lignes sont barrées de cinq longs traits obliques et
d'une rature pleine à droite.}

\emph{Lemme 1}. Soit $A$ un anneau de Dedekind, $M$ un $A$-module plat.
\struck{Conditions équivalentes} Pour $N$ sous-module de $M$, soit
$\widetilde{N}$ l'image inverse \ill{} \ill{} \ill{} de torsion de $M/N$.

\uncertain{D.}\,--- $\widetilde{N}/N$ est un \ill{} de torsion.

Cond.\ équivalentes a) b) c) d)
\note{« a) b) » est écrit deux fois, la seconde fois suivi de « c) d) ».}

a) $\forall N$ de type fini, $\widetilde{N}/N$ est de \uncertain{longueur} finie,
i.e.\ $\widetilde{N}$ est de t.f.
\struck{$A$ est un corps,} si $A$ n'est pas un corps,
\struck{\ill{} \ill{} \ill{} \ill{} \ill{}} \ill{} de $M$
\marginal{un double trait vertical fléché, $\Updownarrow$, en face de a)}

b) \struck{Il n'existe pas de \ill{} \ill{} de
$A(\mathfrak{p}) = \bigcup_n \mathfrak{p}^{-n}
= \Gamma(S - \lbrace\mathfrak{p}\rbrace, \underline{O}_S)$
($S = \mathrm{Spec}\,A$), $\mathfrak{p}$ est un idéal maximal}
\ill{} :
\struck{\ill{} $A$ est un corps}
de $A$
\note{le passage de b) est enfermé dans un cadre et barré ; « $A(\mathfrak{p})$ »
est récrit une seconde fois, surchargé, sous la ligne.}

\marginal{$\Updownarrow$}
Si $A$ n'est pas un corps \ill{} \ill{},
\struck{\ill{} \ill{}}
c) (Pour \ill{} idéal maximal $\mathfrak{p}$ de $A$, $M_{\mathfrak{p}}$
\ill{} \ill{} \ill{} \uncertain{isomorphe} à $K$.

d) $M = \varinjlim M_i$, les $M_i \subset M$ de type fini et « saturés », i.e.\
$M_i \to M_j$ tels que $M_i$ facteur direct dans $M_j$, les applications
\uncertain{de transition} \uncertain{injectives}.
\note{« de type fini et "saturés", i.e.\ \ldots\ facteur direct dans $M_j$ » est
ajouté en interligne au-dessus d'un mot biffé ; la lecture de la suite est
incertaine.}

Pour que ces conditions soient satisfaites, il suffit que l'on ait :

e) les formes linéaires sur $M$ séparent, i.e.\ \uncertain{équivalent}
\[
  M \hookrightarrow A^{I} \qquad (I \text{ un ensemble convenable})
\]
[condition \uncertain{moins fine} : $M$ projectif] ; \ill{} \ill{} $M$ \ill{}
\ill{} \ill{}, \struck{\ill{}} i.e.\ $M$ \ill{} \ill{} \ill{} (c'est aussi
nécessaire si $M$ a un ens.\ dénombrable de générateurs).

\page{2}

\note{feuillet dactylographié, au verso des notes et sans rapport avec
elles ; second feuillet de « EGA IV 18.9. (Rabiots) » (voir page 4), texte
publié, non recomposé. Il porte la fin d'une démonstration de 18.9.7 :
réduction au cas $Y$ géométriquement unibranche (remarque que 19.9.7 est
essentiellement contenu dans son corollaire 18.9.11), puis, via 15.5.6.1, au
cas $\dim Y = 1$ et à la non-déconnexion de $X$ par son point fermé, par
$\mathrm{prof}\,B = \mathrm{prof}\,A + \mathrm{prof}\,B \otimes_A k$ et le
théorème de Hartshorne ; enfin le cas $A$ universellement \ldots, qui ramène
à $A$ normal et à un anneau de valuation discrète. Les flèches et les signes
$\geqslant$, laissés en blanc par la machine, n'ont pas été portés.
Une seule intervention manuscrite : « caténaire » est biffé et remplacé en
interligne par \uncertain{japonais}.}

\page{3}

\emph{Lemme 2}. Soit \struck{\ill{}} $A$ comme dessus, $G$ un $A$-groupe affine,
d'algèbre $B$, satisfaisant aux conditions ci-dessus. Alors pour tout
\struck{\ill{}} sous-module de type fini $M$ de $B$, $\exists$ sous-module
$N$ de $B$ tel que $N \supset M$ et tel que
\begin{itemize}
  \item[a)] $\widetilde{N} = N$ i.e.\ $N$ saturé \ill{} ---
    ce qui implique $N \otimes_A N \hookrightarrow B \otimes_A B$ ;
  \item[b)] $\Delta N \subset N \otimes N$.
\end{itemize}
\note{dans b), le produit tensoriel est d'abord écrit par une boucle
surchargée, puis récrit.}

\emph{Dém.} Faisons opérer $G \times G$ \struck{\ill{}} sur $B$ par
\[
  (g,h)f(s) = f(g^{-1}sh) .
\]
\struck{Alors les} \ill{} \ill{} \ill{} $N$ de $B$ \struck{\ill{}}
\ill{} \ill{} par $G \times e$ (resp.\ $e \times G$) \struck{\ill{}} \ill{}
\ill{} :
\[
  \Delta N \subset N \otimes B \qquad \text{resp.}\qquad
  \Delta N \subset B \otimes N .
\]
Notons que \struck{\ill{} \ill{} la stabilité de $N$, \ill{}} \ill{} \ill{}
\ill{} les conditions \ill{} \ill{} \ill{}
passage de $N$ à $\widetilde{N}$ (\ill{} \ill{}).
\ill{}, la condition d'être invariant par $G \times G$, \ill{} \ill{},
\ill{} \ill{} \ill{} \ill{} \ill{} \ill{}.
N.B.\ \ill{} le passage de $N$ à $\widetilde{N}$ \ill{} \ill{}. \ill{} \ill{}
\ill{} \ill{}\,:
$M \subset N$ \ill{} \ill{} par $G \times G$, \ill{} \ill{} $\widetilde{N}$
\ill{} \ill{} \ill{} de t.f., \struck{\ill{}} \ill{} \ill{} \ill{}.
\marginal{\ill{} \ill{} \ill{} $N = \widetilde{N}$ \ill{}}
Mais cela signifie
\[
  \Delta\widetilde{N} \subset (\widetilde{N} \otimes B) \cap (B \otimes N)
  = N \otimes N
\]
cqfd.
\note{les tildes sont portés sur les deux premiers $N$ de la formule et
manquent aux trois derniers ; transcrit tel quel. L'insertion marginale,
reliée par un trait au-dessus de la dernière ligne, se termine par « $N =
\widetilde{N}$ ».}

\page{4}

\note{feuillet dactylographié, au verso des notes et sans rapport avec elles,
intitulé « EGA IV 18.9. (Rabiots) » ; texte publié, non recomposé. Il donne
le Théorème 18.9.7 (pour $A \to B$ homomorphisme local d'anneaux locaux
noethériens, $\mathrm{Spec}(B) \to \mathrm{Spec}(A)$ plat à fibres séparables,
$A$ géométriquement unibranche ou unibranche avec $k(B)$ extension primaire de
$k(A)$ : la fibre générique est connexe, et $X - T$ est connexe pour tout fermé
$T \subset f^{-1}(S)$, $S \neq Y$ fermé), les Corollaires 18.9.8 (toutes les
fibres géométriquement connexes si $A$ est strictement local, ou hensélien avec
$k(B)$ primaire) avec sa démonstration, et 18.9.9 (fibres formelles connexes
d'un anneau local noethérien hensélien universellement caténaire). Les flèches
sont laissées en blanc. Une seule intervention manuscrite : « géom. » inséré en
marge avant « connexes » à la dernière ligne de 18.9.9.}

\page{5}

\uncertain{Cor. Main}\note{ce titre, souligné, revient page 9 ; il peut se lire
« Corollaire ».}
Si \add{de plus} $G$ est de t.f.\ sur $A$, alors $G$ admet une représentation
linéaire fidèle.
\note{« de plus » est ajouté en interligne ; un trait vertical en marge
gauche tient l'énoncé.}

Il existe $N$ tel que a) $N$ engendre l'algèbre $B$, b) $\widetilde{N} = N$,
c) $\Delta N \subset N \otimes B$. \struck{\ill{}} Alors \struck{\ill{}}
\ill{} \ill{} \ill{} de $G$ : \ill{} \ill{} $A'/A$ qui opère trivialement
dans $N'$, opère trivialement sur $B'$, \struck{\ill{}} donc est l'unité. Donc
la repr.\ est fidèle.

\uncertain{Cor. Main} Si $A$ est un corps $k$, alors \ill{} \ill{} \ill{} $G$
sur $k$ est \ill{} \ill{} \struck{\ill{}} projective d'un syst.\ projectif
strict (\ill{} \ill{} \ill{} morphismes \uncertain{fid.\ plats}) de groupes
affines \emph{de t.f.}

\page{6}

\note{feuillet dactylographié paginé « -3- », au verso des notes et sans
rapport avec elles ; texte publié (EGA IV, 18.9.5), non recomposé. Il porte le
Corollaire 18.9.5 (sous les hypothèses de 18.9.4, le foncteur
$Z \mapsto Z \otimes_A A'$ des préschémas étales sur $X$ dans ceux sur $X'$ est
pleinement fidèle) et sa démonstration, d'abord pour $Z_2$ séparé, puis dans le
cas général par le graphe $\Gamma'$ d'un $X'$-morphisme $Z'_1 \to Z'_2$ ; c'est
une rédaction plus développée que celle, barrée, de la page 8, et elle se
poursuit page 14. Les interventions à l'encre, de main non établie :
\begin{itemize}
  \item « d'un morphisme unique » devient « d'un $X$-morphisme unique » ;
  \item « Supposons d'abord que $Z_1$ et $Z_2$ soient séparés » : « $Z_1$ et »
    est biffé et l'accord ramené au singulier (« $Z_2$ soit séparé ») ;
  \item après « le graphe $\Gamma'$ d'un $X'$-morphisme $Z'_1 \to Z'_2$ »,
    inséré « qui est ouvert (17.9.3) dans $Z'_1 \times_{X'} Z'_2$ » ;
  \item « par changement de base » reçoit en interligne « $A \to A'$ » ;
  \item « la conclusion résulte donc de » devient « \ldots\ donc alors de » ;
  \item « où $\Gamma$ est ouvert » reçoit « dans $Z$ » ;
  \item « un morphisme plat » devient « un morphisme fidèlement plat et
    quasi-compact » ;
  \item deux traits en marge gauche, en face des deux passages corrigés.
\end{itemize}}

\page{7}

\emph{Prop.} \ill{} $A$ \ill{}, $G$ groupe affine sur $A$, d'algèbre $B$. On
suppose que on \ill{} \ill{} \ill{} \ill{} $A$-module,
\[
  B = \varinjlim B_i , \quad \text{où } B_i \subset B \text{ sont tels que}
\]
a) $B/B_i$ est plat$/A$, \struck{\ill{}} donc
$B_i \otimes B_i \subset B \otimes B$, et b) $\Delta B_i \subset B_i \otimes B_i$.
\struck{\ill{}} c) les $B_i$ sont projectifs de t.f.

\struck{Alors} Pour $i$ fixé, utilisant a), on voit que pour tout $A$-module
$E$, $B_i \otimes E \subset B \otimes E$. Si $G$ opère sur $E$, disons que $E$
est \emph{$i$-admissible} si $\delta \colon E \to B \otimes E$ se factorise par
$B_i \otimes E$, donc on définit par
\[
  \delta_i \colon E \to B_i \otimes E .
\]
Alors la catégorie des $A$-$G$-modules \add{de t.f.} est limite inductive,
pour $i$ variable, des sous-catégories $C_i(A,G)_f$ des $A$-$G$-modules
$i$-admissibles \add{de t.f.}
\note{les deux « de t.f. » sont ajoutés en interligne ; dans $C_i(A,G)_f$ la
première lettre de l'argument est surchargée.}
D'autre part, utilisant b) et c), et introduisant $U_i = \check{B}_i$, on
trouve que $U_i$ est une algèbre ass.\ unitaire \struck{\ill{}}
\struck{\ill{} \ill{} \ill{} ainsi que pour que $\varepsilon(B_i) = A$,}
\struck{\ill{} \ill{} \ill{}}
Alors, la catégorie $C_i(U_i)$ \ill{}
\note{les deux avant-dernières lignes sont barrées et encadrées ; la dernière
est barrée à partir de « la catégorie ».}

\page{8}

\note{feuillet dactylographié paginé à la main « IV-965 », au verso des notes
et sans rapport avec elles ; texte publié, non recomposé. Il porte le Corollaire
18.9.4 (pour $A$ local noethérien hensélien à fibres formelles géométriquement
normales, $A' = \widehat{A}$, $g \colon X' \to X$ : $U \mapsto g^{-1}(U)$ est une
bijection des parties ouvertes et fermées, d'où $\pi_0(X') \to \pi_0(X)$
bijectif) et sa démonstration, puis une première rédaction de 18.9.5, barrée de
deux longs traits obliques (la rédaction reprise est celle des pages 6 et 14),
et le début de la Remarque 18.9.6. Les interventions, de main non établie :
\begin{itemize}
  \item la référence à Bourbaki, \emph{Top.\ gén.}, chap.~I, §~11, n°~3, est
    complétée en « prop.~7 » ;
  \item dans l'énoncé barré de 18.9.5 : « $X$-schémas » devient
    « \uncertain{pré}schémas », « (supposés étales et finis) » est ajouté puis
    biffé deux fois, « (et séparés) » ajouté deux fois dans des cadres ;
  \item en marge gauche : « « séparés » inutile ! » et « le vrai
    \uncertain{énoncé} \uncertain{porte} \ill{} \uncertain{sur} $X$-schémas
    \ill{} \ill{} \ill{} \ill{} \ill{}\ldots » ;
  \item dans 18.9.6, « Dans tous les cas connus » est biffé et remplacé par
    « Il est possible que si le corps résiduel de $A$ est de car.\ nulle, » ;
    « est » est remplacé par « induit ».
\end{itemize}}

\page{9}

formant un syst.\ projectif \add{strict} d'algèbres sur $A$.
$C_i(G)_f$ est équivalente canoniquement à la catégorie des $U_i$-modules
\add{de t.f.} Et les foncteurs de transition entre les $C_i(G)$ sont
isomorphes aux \struck{\ill{}} foncteurs \emph{restriction des scalaires}.
\note{« strict » et « de t.f. » sont ajoutés en interligne ; la page suit
directement la page 7.}

\uncertain{Cor. Main} Sous les conditions précédentes, tout
\uncertain{module} \add{de t.f.} $E$ d'une représentation de $G$ est quotient
d'un module \add{projectif} de t.f.\ où $G$ opère.
\note{« de t.f. $E$ » est ajouté en interligne avec un signe de renvoi ;
« projectif », souligné, est ajouté au-dessus d'un mot biffé.}

\emph{Prop.} Soit $G$ groupe affine plat sur $A$ anneau de Dedekind, d'algèbre
$B$ satisfaisant aux conditions du Lemme 1 (cela signifie que pour les
composantes irréd.\ de $G$, \emph{leurs} fibres sont non vides).
Soit \struck{\ill{} \ill{} \ill{}} $S'$ un \uncertain{schéma} sur $S$, \ill{}
\ill{} sur $S$, \ill{} module quasi-coh.\ \struck{\ill{}} \uncertain{soit}
de t.f.\ quotient d'un module \uncertain{localement libre} de t.f.
\struck{Alors} (\uncertain{p.\ ex.} $S$ \ill{}, \ill{} $S$ noeth.\
régulier). Alors \ill{}

\page{10}

\note{feuillet dactylographié paginé à la main « IV-964 », au verso des notes
et sans rapport avec elles ; texte publié, non recomposé. Il porte la fin d'un
énoncé sur les anneaux unibranches, la Proposition 18.8.15 ($^{hs}A$
universellement caténaire si $A$ l'est), le début du §~18.9 « Fibres formelles
des anneaux noethériens henséliens » : Théorème 18.9.1 (fibres formelles
géométriquement intègres), Corollaire 18.9.2 ($\widehat{A}$ intègre) et
Corollaire 18.9.3 (le corps des fractions de $\widehat{A}$ est une extension
séparable de celui de $A$, dans laquelle ce dernier est algébriquement fermé),
avec leurs démonstrations. Les interventions, de main non établie :
\begin{itemize}
  \item le numéro « (18.8.13) » est corrigé en « (18.8.15) », « (18.?.3) » en
    « (18.9.3) » ;
  \item « et » ajouté devant « $K$ est algébriquement fermé dans $L$ » ;
  \item « La première assertion » est biffé et remplacé par « Cela » ; un mot
    biffé devant « intègre » ;
  \item la référence finale est complétée : « et de (4.3.2) et (4.3.5) ».
\end{itemize}}

\page{11}

$G_{S'}$-modules \struck{\ill{}} quasi-coh.\ de t.f.\ est quotient d'un idem,
qui est de plus \emph{loc.\ libre}.
\note{suite directe de la page 9 ; le reste du feuillet est blanc.}

\page{12}

\note{feuillet dactylographié paginé à la main « IV-9\ill{} », au verso des
notes et sans rapport avec elles ; texte publié, non recomposé. Il porte la fin
d'un énoncé sur $^{hs}A$ (fibres géométriquement régulières), les Corollaires
18.8.12 (Cohen-Macaulay, $(S_n)$, régulier, $(R_n)$ passent de $A$ à $^{hs}A$
et inversement) et 18.8.13 (idéaux premiers associés), la Proposition 18.8.14
(équivalence de $A$ géométriquement unibranche, de l'irréductibilité de
$\mathrm{Spec}(A')$ pour $A'$ essentiellement étale, et de celle de
$\mathrm{Spec}(^{hs}A)$) avec sa démonstration, et le début du Corollaire
18.8.15. Les interventions, de main non établie :
\begin{itemize}
  \item les numéros sont décalés à l'encre : 18.8.9, 10, 11, 12 deviennent
    18.8.12, 13, 14, 15 ; dans le dernier, un mot tapé est surchargé en
    « Corollaire » ;
  \item une accolade en marge réunit b), c), d) de 18.8.12, avec un signe
    \ill{} ;
  \item « en utilisant (18.8.7.1) » est mis entre crochets et souligné d'un
    trait ondulé ; un trait en marge ;
  \item dans 18.8.15, « local » est inséré devant « hensélien ».
\end{itemize}}

\page{13}

\struck{\emph{Proposition} Soit $A$ anneau de Dedekind, $G$ gr.\
\add{plat} affine \add{de t.f.} sur $A$, d'algèbre $B$ satisfaisant aux
conditions du Lemme 1. \struck{Alors} Soit $N$ comme dans le Lemme 2, et
considérons la représentation de $G$ dans $N$ induite par la repr.\ régulière
gauche. Alors l'hom.\
$G \to \underline{\mathrm{Aut}}(N)$ \struck{est} \ill{} un \uncertain{morphisme}
\ill{} \uncertain{plat}}
\note{tout l'énoncé est barré de six longs traits obliques et séparé de la
suite par un trait horizontal.}

\[
  \left\lbrace
  \begin{array}{l}
    \mathcal{B} = \varinjlim \mathcal{B}_i \\
    1)\ \mathcal{B}/\mathcal{B}_i \text{ plat, donc }
      \mathcal{B}_i \otimes \mathcal{B}_i \subset \mathcal{B} \otimes \mathcal{B} \\
    2)\ \Delta \mathcal{B}_i \subset \mathcal{B}_i \otimes \mathcal{B}_i \\
    3)\ \mathcal{B}_i \text{ loc. libres de t.f.}
  \end{array}
  \right.
\]
\note{après 3), « Supposons $X$ quasi-proj.\ \ill{} » est biffé.}

Soit $E$ un $\underline{O}_X$-$G$-module quasi-coh. \struck{\ill{}}

Alors 1) Pour tout sous-module quasi-coh.\ $F$, il existe un plus petit
sous-module qu.-coh.\ \ill{} qui contienne $F$, \struck{\ill{} \ill{}}
stable par $G$, et \ill{} \ill{} sa formation commute \ill{} tout changement de
base plat. Si $F$ est de t.f., $F'$ itou.

2) Si $X$ quasi-compact quasi-séparé, alors $E$ est \emph{lim.\ ind.}\ de
sous-modules \add{de t.f.}\ qu.-coh.\ stables par $G$.

3) Soit $F \to E$ un \ill{} \ill{}, \struck{\ill{}} $F$ quasi-coh.\
\add{de t.f.}, \ill{} l'image engendre $E$. Alors il existe
\note{entre 2) et 3), un passage est biffé et enfermé dans un cadre :
« \ill{} \ill{} \ill{} si $X$ quasi-compact et \uncertain{principal} \ill{} ».}

\page{14}

\note{feuillet dactylographié paginé « -4- », au verso des notes et sans
rapport avec elles ; fin de la démonstration du Corollaire 18.9.5 commencée
page 6 (réduction, via 4.5.19.1, à $q_1^{-1}(\Gamma') = q_2^{-1}(\Gamma')$, puis
au cas où $X$ est le spectre d'un corps, où tout $X$-préschéma étale est
séparé, 17.6.2, c')). Seule intervention : la lettre finale de « est »,
complétée à la main.}

\page{15}

un $G$-hom.\ \struck{\ill{}} $F' \to E$ surjectif, avec $F'$ de la forme
$F \otimes \check{\mathcal{B}}_i$ --- donc $F'$ loc.\ libre (de t.f.) si $F$
l'est.
\note{suite directe de la page 13.}

4) Supposons $X$ \uncertain{quasi-proj.}\ \ill{} \uncertain{et} que \ill{}
module \ill{} de t.f.\ sur $X$ \ill{} quotient d'un loc.\ libre de t.f.\
(\uncertain{p.\ ex.} $X$ \uncertain{admet} faisceau inversible ample, ou $X$
noethérien \uncertain{régulier}). Alors \ill{}
$\underline{O}_X$-$G$-module sur $X$ est quotient d'un loc.\ libre de t.f.

\page{16}

\note{feuillet dactylographié paginé à la main « IV-962 », au verso des notes
et sans rapport avec elles ; texte publié, non recomposé. Il porte la fin de
la démonstration de 18.8.7 (pour $C$ fini sur $A$ : $C \otimes_A {}^{hs}\!A$
composé direct d'anneaux strictement locaux $C'_i$, construction de
$w \colon C'_i \to {}^{hs}C$ et de son inverse $w'$), la Remarque 18.8.10
(étendue à $B$ entière semi-locale à extensions résiduelles de degré séparable
fini) et le début de la Proposition 18.8.11 (réduit, normal ; $(^{hs}A)_{\mathfrak
p}/\mathfrak{p}(^{hs}A)_{\mathfrak p}$ composé d'extensions séparables). Les
interventions, de main non établie :
\begin{itemize}
  \item un diagramme tracé à la main, sans lequel le texte ne se suit pas :
    $A \to {}^{hs}\!A$, $A \to C$, $f \colon C \to C'$, $C' \to C'_i$,
    $v \colon C \to {}^{hs}C$, $w \colon C'_i \to {}^{hs}C$,
    $u \colon {}^{hs}\!A \to {}^{hs}C$, avec en marge
    « $C \to {}^{hs}C$ » pour le blanc laissé après « Soit $v$ » ;
  \item les relations $w_0 \circ g = u$, $w_0 \circ f = v$ et la flèche
    $\varphi_i$ sont portées à la main dans des blancs ; « (18.8.5) » ajouté
    au-dessus de $w' \circ v = \varphi_i \circ f$ ;
  \item les numéros « (18.8.7.1) » et « (18.8.8) » deviennent « (18.8.10) »
    et « (18.8.11) » ;
  \item en marge, dans une ellipse au crayon : « pour une $A$-algèbre $B$
    \emph{entière} sur $A$ », relié par un trait au texte de la Remarque.
\end{itemize}}

\page{17}

$G$ un groupe plat et affine sur $S$, d'algèbre $\mathcal{B}$. On suppose
\[
  \mathcal{B} = \varinjlim \mathcal{B}_i , \qquad
  \mathcal{B}_i \subset \mathcal{B} \text{ sous-Modules avec}
\]
\begin{itemize}
  \item[1)] $\mathcal{B}/\mathcal{B}_i$ plat (donc
    $\mathcal{B}_i \otimes \mathcal{B}_i \subset \mathcal{B} \otimes \mathcal{B}$)
  \item[2)] $\Delta \mathcal{B}_i \subset \mathcal{B}_i \otimes \mathcal{B}_i$
    donc $\check{\mathcal{B}}_i$ est un faisceau d'algèbres associatives à unité
  \item[3)] $\mathcal{B}_i$ \struck{projectif} loc.\ libre de t.f.\ donc les
    $\check{\mathcal{B}}_i$ itou.
\end{itemize}
\note{« loc.\ libre » est écrit au-dessus de « projectif », biffé.}

$X$ un préschéma sur $S$, \struck{muni de $\underline{L}$ ample rel$/S$}
\uncertain{où} $G$ opère. \struck{Groupe} On suppose que sur $X$ existe un
faisceau \add{inversible} $\underline{L}$ ample rel$/S$, sur lequel $G$ opère
de façon compatible avec son action sur $X$.

Soit $E$ un $\underline{O}_X$-$G$-Module quasi-coh. \struck{\ill{}}
\struck{\ill{}}

1) Pour tout sous-module quasi-coh.\ $F$ de $E$, il existe un plus petit
sous-module quasi-coh.\ $F'$ de $E$ \struck{qui} stable par $G$ contenant $F$.
Sa \struck{\ill{}} formation commute à tout changement de base plat en $S$.

\struck{2) Si $S$ \ill{} quasi-compact quasi-séparé, alors $E$ est
\emph{lim.} de sous-modules quasi-coh.\ de t.f.\ \emph{stables} par $G$.}
\note{ce passage est encadré et barré de traits obliques ; en marge gauche,
biffé, « \uncertain{de t.f.} ».}

2) Si $F$ est de t.f., de même $F'$.
\note{inséré en interligne au-dessus du passage barré ; le numéro 2) est
porté en tête du passage barré et repris pour l'insertion.}

\page{18}

\note{feuillet dactylographié paginé à la main « IV-953 », au verso des notes
et sans rapport avec elles ; texte publié, non recomposé. Il porte la fin
d'une démonstration (espace compact noethérien, donc fini et discret), le
Corollaire 18.6.10 (Cohen-Macaulay, $(S_n)$, régulier, $(R_n)$ pour $^{h}\!A$),
le Corollaire 18.6.11 (idéaux premiers associés) et la Proposition 18.6.12
($A$ unibranche ssi $\mathrm{Spec}(^{h}\!A)$ irréductible) avec le début de sa
démonstration. Les interventions, de main non établie : une accolade en marge
réunit b), c), d) de 18.6.10 avec un signe \ill{} cerclé ; les chiffres
finaux de « (18.6.9) » et « (18.6.12) » sont surchargés ; un trait en marge de
18.6.12 ; « Soient » corrigé en « Soit ».}

\page{19}

\struck{3) Si $\underline{E}$ est de t.f., il existe un
$\underline{O}_X$-$G$-Module \ill{} \emph{loc.\ libre} \emph{de t.f.}\
$\underline{F}$, et un épimorphisme de $\underline{O}_X$-$G$-Modules
$\underline{F} \to \underline{E}$.}
\note{ce 3) est barré de traits obliques ; au-dessus, ajouté puis biffé avec
lui : « et est quotient d'un module loc.\ libre (sans op.\,!) ».}

3) \struck{Si $S$ \ill{} \ill{} \ill{} \ill{} faisceau \ill{}}
\add{\uncertain{quasi-proj.}\ \ill{} et si tt module quasi-coh.\ de t.f.\ sur
$S$ est quotient d'un loc.\ libre de t.f.,} \struck{\ill{}} alors
$\underline{E}$ est quotient d'un $\underline{O}_X$-$G$-Module loc.\ libre de
type fini de la forme $f^{*}(H)(-n)$, \ill{}
\note{la condition sur $S$ est écrite en interligne au-dessus d'une ligne
biffée.}

\emph{Dém.} Il suffit de prouver 1) et 3). Posons (où $f \colon X \to S$)
\[
  \underline{\Gamma}_{*}(E) = \coprod_n f_{*}(E(n)) .
\]
\note{le signe de somme est un $\coprod$, $n$ en dessous ; la page emploie
ensuite $\underline{\Gamma}^{\cdot}$ pour le même objet.}
Dans \uncertain{le cas} où $G$ est plat sur $S$ et $f$ qu.-cpt et
(qu.-)sép., on trouve que $\underline{\Gamma}^{\cdot}(E)$ est un Module gradué
sur $S$ sur lequel $G$ opère.
\uncertain{Si}
\[
  \underline{\Gamma}^{\cdot}(\underline{O}_X) = \struck{\ill{}}\, A^{\cdot}
\]
c'est une Algèbre graduée sur laquelle $G$ opère. $E \mapsto
\underline{\Gamma}^{\cdot}(E)$ est un foncteur \uncertain{de} $E$, à valeurs
dans la cat.\ des $A^{\cdot}$-$G$-modules gradués. Ce foncteur est exact à
gauche. On a un foncteur (\emph{exact})
\note{la phrase se poursuit en tête de la page 21 (batch 2) : « $M \mapsto
\widetilde{M}$ en sens inverse ».}

\page{20}

\note{feuillet dactylographié paginé à la main « IV-95\ill{} », au verso des
notes et sans rapport avec elles ; texte publié, non recomposé. Il porte la fin
de l'énoncé de 18.6.9 (fibres de $\mathrm{Spec}(^{h}\!A) \to \mathrm{Spec}(A)$
géométriquement régulières), la démonstration de (i) (réduit, normal) et de
(ii), et les Lemmes 18.6.9.1 (limite inductive noethérienne d'anneaux
artiniens) et 18.6.9.2 (limite projective noethérienne d'espaces finis
discrets). Les interventions, de main non établie :
\begin{itemize}
  \item « les morphismes $\mathrm{Spec}(A_\beta) \to \mathrm{Spec}(A_\alpha)$ »
    devient « $A_\alpha \to A_\beta$ » ; « surjectifs » devient « injectifs » ;
    la référence « (2.3.4, (ii)) » est barrée ;
  \item en marge gauche : « \uncertain{précisions} \uncertain{mélangées} ! » ;
  \item « (On peut se borner au cas où $A$ est local.) » est rattaché à (ii)
    par une parenthèse ;
  \item « $\varphi_{\beta\alpha}$ » porté à la main au-dessus de « les
    homomorphismes $B_\alpha \to B_\beta$ », puis dans la démonstration ;
  \item « étant discrets » devient « finis et discrets » ;
  \item une note au crayon en marge gauche, \ill{} \ill{} \ill{}.
\end{itemize}}

\end{document}
